%!TEX root = ../report.tex
\documentclass[report.tex]{subfiles}
\begin{document}
    \chapter{Methodology} \label{chap:methodology}

    This chapter contains information on the how part of the problem statement. What techniques are used to estimate the odometry using InEKF-pcv
        
    \section{Proposed approach}
    
    The proposed algorithm for odometry estimation is based on fusion of \gls{imu} data and the wheel encoder data. For the fusion process, the \gls{pcv} vector is calculated. The \gls{pcv} is calculated in two forms. The \gls{pcv} vector, calculated from the filter state vector after the prediction step of the \gls{inekf} is called predicted PCV. In the update step, the \gls{pcv} vector is derived from the wheel encoder measurement called measured PCV. Within the \gls{inekf}, the predicted and measured \gls{pcv} vectors are used to calculate the residual, also known as innovation.
    
    \section{Invariant Extended Kalman Filter}
    
    The background details relevant to the understanding of the \gls{inekf} is provided in the background chapter \ref{chap:background}. In this section, the reason behind the \gls{inekf} choice for odometry estimation in \gls{rwr} is discussed.
    
    \todo[inline]{need to get a better understanding of why Inekf is better here, and explain here}
    The \gls{inekf} is symmetry precerving observer, meaning the rotation matrix in the state vector is handelled as a whole. 
    
    
    The following subsections will provide details on the state vector, the prediction step, the update step, the measurement function and its corresponding jacobian.
    
    \subsection{Filter State Vector}
    
    The filter state is constructed as shown in the equation \ref{eq:stateVector}. 
    \begin{equation}
        X = [R, V, P, b_a, b_g]
        \label{eq:stateVector}
    \end{equation}
    Where, $R$ is the Rotation matrix, that represents the rotation of body in world. $V$ is the velocity of body in world, $P$ is the position of body in world, $b_a$ is the acceleration bias, $b_g$ is the gyro bias. 
    
    \subsection{Propagation step} \label{subsec:inekfPropagationStep}
    The \gls{imu} data is used to propagate the state from current time $t$ to the next time step $t+\delta t$. The filter state is propagated in time using a state transition matrix and a noise matrix. \todo[inline]{This is very vague, elaborate after adding the content to the background} The propagation step in the \gls{inekf} is left unchanged, the original implementation provided in the \gls{inekf} library \cite{REPO-inekf-github}.
    The propagation step is out of scope for this thesis and the existing implementation is used. 
    
    \subsection{Update step} \label{subsec:inekfUpdateStep}

    The update step is the focus of this thesis, where the predicted \gls{pcv} and measured \gls{pcv} are calculated. This subsection explains the steps to calculate the predicted and measured \gls{pcv}.
    
    \subsubsection{Predicted \gls{pcv}}
    
    The predicted \gls{pcv} can be calculated from the state vector \ref{eq:stateVector} after the \gls{inekf} propagation step. All the calculations related to the \gls{inekf} are body centric calculations. \todo[inline]{What does body centric calculation mean? Elaborate here} In order to calculate the velocity of a point in a rigid body relative to the body frame, we need to represent the velocity of body frame at the point we need the velocity at. In our case, our robot model has a body frame, which is located at the center of the rover, and we need to calculate the velocity of the \gls{pc} relative to the body frame. 
    
    This calculation can be done using the equation \ref{eq:ExpectedPCVCalculation}. This operation can also be done using the \gls{kdl} library using Twists \ref{subsec:twists} by utilizing the adjoint representation of the transformation matrix using the equation \ref{eq:TwistTransformationEquation}. The implemented version of the twist transformation is provided in the equation \ref{eq:ExpectedPCVCalculationTwist}
    
        \begin{equation}
            v_{pc}^{body} = v_{body}^{body} + (\omega_{body}^{body} \times {}^{}r_{pc}^{body})
            \label{eq:ExpectedPCVCalculation}
        \end{equation}
                
        \begin{equation}
            \begin{bmatrix}
            \boldsymbol{\omega}_{pc}^{body} \\
            \mathbf{v}_{pc}^{body}
            \end{bmatrix}
            =
            \begin{bmatrix}
            I & 0 \\
            \left[\,\mathbf{p}_{pc}^{body}\,\right]_\times & I
            \end{bmatrix}
            \begin{bmatrix}
            \boldsymbol{\omega}_{body}^{body} \\
            \mathbf{v}_{body}^{body}
            \end{bmatrix}
            \label{eq:ExpectedPCVCalculationTwist}            
        \end{equation}
        
    % The calculated $v_{pc}^{body}$ represents the velocity of the point(s) that is in contact with the ground.
    
    \subsubsection{Measured \gls{pcv}}        

        
    \subsection{Measured Potential Contact Velocity Calculation}
        
    \subsection{Fixed Wheel Frames}
    
    The fixed wheel frames are added in the robot URDF model at the location of motors. This fixed wheel frames have the same pose as the initial pose of the rotating wheel frame. The fixed wheel frames are added to get the position vector to the spoke of the wheel from the center of the wheel. These position vectors are later used to calculate the measured \gls{pcv}. This is by performing a cross product of wheel's angular velocity vector (the vector form of rotation speed) and the position vector of \gls{pc}. The \gls{pcv} vector will be perpendicular to the spoke and point in the direction of rotation of the wheel. This is the only purpose of the fixed wheel frames. The calculation of the measured \gls{pcv} is shown in equation \ref{eq:MeasuredPCVCalculation}.
    
    \begin{equation}
        v_{pc}^{wheel} = \omega_{wheel}^{wheel} \times {}^{}r_{pc}^{wheel}
        \label{eq:MeasuredPCVCalculation}
    \end{equation}
    
    
    \subsection{Measurement function} \label{subsec:inekfMeasurementFunction}
    
    \subsection{Measurement Jacobian} \label{subsec:inekfMeasurementJacobian}
    
    
    \subsection{Measurement Jacobian}
    
        The measurement jacobian $H$ is the partial derivative of the measurement function $h$ with respect to the error state $\xi$ variables at the state $\hat{X}$ as shown in equation \ref{eq:MeasurementJacobian-1}. The measurement jacobian linearizes the measurement function at the current state. The measurement function $h(x)$ is given by the equation \ref{eq:MeasurementFunction}. Measurement jacobian for this measurement function is obtained by taking the partial derivative of the measurement function with respect to the error state variables \ref{eq:ErrorStateVector}. The measurement jacobian calculation and its results for this measurement function is shown in equations \ref{eq:MeasurementJacobian-2} and \ref{eq:MeasurementJacobian-3} respectively.
    
        \begin{equation}
            h(x) = - (R_{world}^{body} * v_{body}^{world} + (\omega_{body}^{body} \times {}^{}r_{pc}^{body}))
            \label{eq:MeasurementFunction}
        \end{equation}
        
        \begin{equation}
            \xi = \begin{bmatrix}
            \delta R \\
            \delta V \\
            \delta P \\
            \delta acc\_bias \\
            \delta gyro\_bias
            \end{bmatrix}
            \label{eq:ErrorStateVector}
        \end{equation}
        
        \begin{equation}
            H = \frac{\partial h}{\partial \xi}\bigg|_{\hat{X}}
            \label{eq:MeasurementJacobian-1}
        \end{equation}
        where, $\xi$ is the left invariant error state, $\hat{X}$ is the current state estimate, and $h$ is the measurement function. 
        
        
        \begin{equation}
            H = \begin{bmatrix}
            \frac{\partial h}{\partial \delta R} & \frac{\partial h}{\partial \delta V} & \frac{\partial h}{\partial \delta P} & \frac{\partial h}{\partial \delta acc\_bias} & \frac{\partial h}{\partial \delta gyro\_bias}
            \end{bmatrix}
            \label{eq:MeasurementJacobian-2}
        \end{equation}
        
        \begin{equation}
            H = \begin{bmatrix}
                -[R_{world}^{body} * v_{body}^{world}]_\times & -\mathbf{I} & \mathbf{0} & \mathbf{0} & -[r_{pc}^{body}]_\times
            \end{bmatrix}
            \label{eq:MeasurementJacobian-3}      
        \end{equation}
       
        The evaluation of equation \ref{eq:MeasurementJacobian-2} to get the result in equation \ref{eq:MeasurementJacobian-3} is discussed in detail in the appendix \ref{appendix:MeasurementJacobianDerivation}.

    \subsection{Slip elimination}
    
    The slippage is discussed in detail in the background chapter \ref{subsec:backgroundSlippage}
    
    \subsection{Contact estimation}
    
    The design of rimless wheel, being a hybrid between wheeled and legged systems have the swing phase and the contact phase. Estimating which spoke is currently is in contact for a given wheel is a vital information for the Contact Velocity Odometry estimation. There is a traditional way of estimating the contact spoke by using the assumption that the spoke that is close to the ground is in contact. This can be calculated by calculating the robot kinematics, using the joint positions of the wheel. The pose of each spoke with respect to the body frame can be calculated and the spoke in contact can be estimated. This assumption works well in a flat, solid terrain but on an inclined or rough terrain with obstacles, this assumption does not hold. 
    
    The use of Kalman filter for odometry estimation provides an opportunity to use its framework to estimate the feet that are in contact.
    In order to understand the contact estimation process, one must understand the Kalman filter framework implemented in this thesis and explained in \ref{subsec:inekf_implementation} \todo[inline]{Write a subsection on InEKF implementation}. The Kalman filter calculates the term innovation, in the case of this thesis, the innovation is calculated as the difference between the velocity of spoke from the filter's prediction and the velocity of the spoke 
    
    

    
\end{document}
